\section{Discussion}
\label{sec:discussion}

The combined evidence supports a positive but conditional view of adaptive
context selection. State-conditioned marginal utility is a useful decision
object, predictor response is the observable that exposes its candidate-specific
variation, and response-aware selection can substantially outperform content,
fixed-set and response-free alternatives when that variation is identifiable.
The same theory also explains why this advantage is not monotone in predictor
quality or candidate-library size.

\subsection{Structural opportunity and realizable opportunity}

Structural headroom measures what an oracle could gain by revising candidate
values after each addition. Realizable opportunity measures how much of that
gain can be predicted from information available before the outcome is known.
The two quantities need not track one another. In the linear identifiable
regime, large oracle opportunity coincides with measurable candidate-specific
signal and response-aware routing converts part of it into gain. With the
strong nonlinear expert, the oracle opportunity remains large while candidate
rankability collapses: cached recall approaches its random reference, observable
probes have near-zero out-of-sample predictive value, and episode-level
confidence signals have AUCs in $[.481,.521]$. The loss is therefore not
explained by a missing scoring head alone.

This distinction also clarifies the expert-strength and library-size sweeps.
Adapting the predictor removes systematic residual structure and sharply reduces
the value a router can realise even when the oracle opportunity persists.
Increasing the library from 16 to 128 candidates raises the oracle opportunity
by about 40\% while the router approaches random selection. More candidate
choice can therefore increase the value of perfect information and
simultaneously reduce the fraction of that value recoverable from observable
signals.

\subsection{Why response matters}

The theory and the empirical probes point to the same interaction. Under squared
loss, the state-dependent part of utility is the alignment between the
candidate-induced response and the residual of the current predictor. A
content-only score omits that interaction. The exact Bregman identity and the
cross-entropy bound show that the same geometric structure is not peculiar to
squared error. The theory-shaped head improves ranking where the signal exists,
while four alternative parameterisations hit the same ceiling in the
low-identifiability regimes. This is why we treat response as a principled
observation rather than as another feature engineering choice.

\subsection{Pool geometry is actionable}

The pool-geometry sweep reveals a complementary mechanism. Relevance is harmful
when highly correlated candidates contribute little information not already
present in the anchor, and becomes useful as candidate novelty increases. The
relationship is strong within the controlled deployment geometries but does not
support a universal cross-dataset threshold. The resulting adaptive rule is
therefore deliberately local: estimate novelty on the deployment's training
split, choose the threshold on held-out targets, and use it to decide whether a
retrieval-style rule is appropriate or whether pooling is the safer default.

This mechanism is useful even when it does not beat a learned router. It turns a
selection failure into a deployment choice and provides a label-free fallback
for settings where learning state-conditioned utility is not justified.

\subsection{Ranking, stopping and task family}

The demonstration-selection family exposes a different limit. There the
standalone and sequential oracles are nearly identical under the natural
stopping rule, so dynamic ranking has little structural value to recover.
Forcing the budget restores a consistent advantage from state-conditioned
ranking, separating the ranking problem from calibration of the stopping rule.
This matters because a selector can fail either because candidate ordering is
unidentifiable or because the deployment rule removes the opportunity before
ranking matters. These are different regimes and require different remedies.

\subsection{Practical deployment}

The final output of the study is a procedure rather than a single unconditional
policy. A deployment should first establish that the candidate pool offers
structural opportunity, then test whether candidate-level utility is
predictable from the available state and response, and finally choose between
response-aware routing, a simple geometry-adaptive rule and pooling. The pilot
experiments show that this decision can be made with far fewer labelled targets
than a full deployment audit, while the time-shift analysis warns that
confidence transfers less reliably than direction. Split-half consistency is
therefore a certification device for strong pilot decisions, not a substitute
for additional labels on marginal cases.

\subsection{Limitations}

The general-loss analysis characterises the information carried by predictor
response, but it does not make utility label-free: the state-dependent
coefficient still has to be learned from outcomes. The strong-backbone study
uses a finite set of architectures and the second task family uses frozen CLIP
representations, so other inductive biases may expose different observable
signals. The novelty rule is deployment-specific and its threshold should not
be transferred blindly across datasets. Finally, the one-step regret statement
does not provide a global guarantee for arbitrary non-monotone set utilities.
These limitations define what must be re-measured when the predictor, candidate
pool or task changes; they do not alter the central distinction between
structural and identifiable opportunity.
